\cotefile{52}
\scanpage{2}

Th. de Torelli

% archive sheet number: 1

\scanpage{3}

\uline{Th. de Torelli}

\uline{\unsure{Th)1|Th 1|Th. 1}} Soit $C$ un schéma en courbes simple, propre, irréductible sur $S$ \inserted{loc. \unsure{noeth|noeth.}}, de genre $g\geq 2$, $J^0_{C/S}=\underline{\underline{\mathrm{Pic}}}^0_{C/S}$ son schéma jacobien homogène. Alors $\mathrm{Aut}_S(C)\longrightarrow\mathrm{Aut}_{S\text{-gr.}}(J^0_{C/S})$ est \uline{injectif}

Soit $\sigma$ un automorphisme de $C$ induisant l'identité sur $J^0_{C/S}$. Soit $\bar{\sigma}$ l'automorphisme de $J^{\unsure{x|\times}}_{C/S}$ induit par $\sigma$, il sera donc de la forme $T_s$, translation associée à une section $s$ de $J^0_{C/S}$ (lemme à énoncer, \deleted{d'ailleurs trivial car se} \inserted{qui moyennant dualité devrait se} ramener à l'énoncé ensembliste). Il faut montrer que la section $s$ est l'identité. À vrai dire, $C$ étant projectif sur $S$, l'énoncé est équivalent au
% slanted
\margin{\uline{\unsure{dégager|dispos|digres}} \unsure{?}}

\uline{Corollaire 1} $\underline{\underline{\mathrm{Aut}}}_S(C)\longrightarrow\underline{\underline{\mathrm{Aut}}}_{S\text{-gr}}(J^0_{C/S})$ est un monomorphisme (et même une \uline{immersion fermée}, si on se rappelle que $\underline{\underline{\mathrm{Aut}}}_S(C)$ est \uline{propre} sur $S$)
% slanted
\margin{permuter le corollaire et le th.}
\deleted{\unsure{que} $\underline{\underline{\mathrm{Aut}}}_S(C)$ est \unsure{la fibre unité} de $\underline{\underline{\mathrm{Aut}}}_{S\text{-gr}}(\unsure{J^0_{C/S}})$ \unsure{et devient} \illegible{2}}

Comme les deux mbs sont \uline{non ramifiés} sur $S$ [pas de automorphismes infinitésimaux !] on est ramené \deleted{\unsure{au}} à prouver que c'est injectif \uline{géométriquement} (\deleted{\unsure{fibre}} \unsure{géom.} par fibre géom., et même pts : \unsure{schéma de base un|allons donc dire un|valeurs dans des} corps alg. clos), donc à énoncer le

% archive sheet number: 2

\scanpage{4}

\uline{Lemme} Soit $C$ une courbe alg. non sing. complète \deleted{\unsure{alg|et}} irréd. de genre $g\geq 2$ sur $k$ alg. clos, $\sigma$ un automorphisme de $C$ induisant l'identité sur $J^0_{C/S}$. Alors $\sigma$ est l'identité.

En effet, le morphisme sur $J^1_{\deleted{C/S}}$ induit par $s$ est de la forme $T_a$, où $a$ est un pt de $J^0$ \unsure{rat|rat.} $/k$, égal à $\sigma x-x$ pour tt $x\in C(k)$. ($C$ étant considéré comme une \unsure{ptie|partie|plongé} de $J^1$).

Donc $\sigma(x)=x+a$, donc ce lemme équivaut au suivant (ds le cas particulier $S=\mathrm{Spec}(k)$)

\uline{Corollaire 2} \uline{Sous les conditions du th}, \deleted{Toute} \inserted{Une} section $s$ de $J^0_{C/S}$ sur $S$ telle que la translation $T_s$ de $J^1_{C/S}$ invarie $C$ (plongée dans $J^1_{C/S}$) est \deleted{l'identité} \inserted{section nulle}

En effet, il résulte de la formule des Traces que le ``\unsure{nb|nb.} des pts. fixes'' \deleted{de} $\sigma$ est égal à
\[
\deleted{1-2g+1=2-2g}\quad \inserted{-1+2g-1=2(g-1)>0}
\]
(puisque $\sigma$ opère trivialement dans la jacobienne $J^0$, donc sa trace y est $2g$] et si $x$ est un pt fixe, on a $a=\sigma x-x=0$

% archive sheet number: 3

\scanpage{5}

\deleted{Utilisant le corollaire 2 \deleted{lorsque $S=\mathrm{Spec}(k\unsure{[\epsilon]|[t]|(\varepsilon)})$} pour $S'=\mathrm{Spec}(\underline D(\underline{\unsure{I|J}}))$ où $\underline I$ est un faisceau coh. sur $S$, $\underline D(\underline I)=\underline{\mathcal O}_S+\underline I$ ($\underline I$ de carré nul) et tous les $S'$-automorphismes de $C'=C\times_S S'$ qui induisent l'identité sur $S$, on trouve}

\deleted{\unsure{Prop.}} \deleted{\unsure{Partons}} \uline{Corollaire 3} Soient $C,C'$ deux courbes sur $S$ satisfaisant les hyp. du th. Alors
\[
\underline{\underline{\mathrm{Is}}}_S(C,C')\longrightarrow\underline{\underline{\mathrm{Is}}}_S(J^0_{C/S},J^0_{C'/S})\quad\text{est une}
\]
immersion fermée.

Soit maintenant $A$ un \uline{schéma abélien} \uline{polarisé} sur $S$ \inserted{de dim relative $g$} (polarisé au sens strict : on se donne une section du schéma de \unsure{Néron-Sévéri|Néron-Séveri|Nér.-Sév.}). Pour \unsure{tt} $\deleted{S}T'$ sur $S$, on désigne par $F(T)$ l'ensemble des classes, à isom près, de courbes $C$ sur $\deleted{\unsure{S'|S_T}}T'$, satisfaisant aux conditions générales de ce \unsure{N°|Th.}, et \uline{munies} d'un isom de $J^0_{C/T}$ sur $A$ \deleted{\unsure{tel de façon que}} \inserted{compatible avec} les polarisations.

% archive sheet number: 4

\scanpage{6}

\deleted{$F'(T)$ l'ensemble analogue, mais où dans la structure de $C$ on se donne de plus une trivialisation \deleted{\unsure{par}} de \deleted{$\unsure{\Omega^1|\underline{\Omega}}$} l'image directe de $(\Omega^1_{C/T})^{\otimes\unsure{h|k}}$ ($h\geq 3$ étant choisi). Soit
\[
\unsure{\nu|d|v}=2h(g-1)-(g-1)=(2h-1)(g-1)
\]}

\deleted{Proposition} \inserted{\uline{Théorème}} \unsure{Les|Le} foncteurs $F(T)$ \deleted{et $\unsure{F'(T)|F(T)}$ sont} \inserted{est} représentable, \deleted{et} en un $M$ quasi projectif sur $S$ ; le morphisme \deleted{c'est un fibré principal homogène en $F(T)$ \unsure{est} \inserted{représenté \unsure{par}} \uline{un sous-schéma} (\unsure{au} \unsure{quasi} \unsure{non fermé}) de $S$, et $F'(T)$ est un fibré principal homogène sur $F(T)$, de groupe $\underline{\underline{\mathrm{GL}}}(\unsure{\nu})_S$} \inserted{structural $M\to S$ est radiciel et \unsure{a une|à une} \unsure{image} \deleted{\unsure{relativement|rel.}} fermée.}
% slanted
\margin{\unsure{rerédiger|rédiger} proprement}

Plus précisément : le morphisme structural \deleted{\unsure{est}} \inserted{\unsure{induit}} \deleted{\unsure{un}} \inserted{un} morphisme \deleted{\unsure{propre et quasi-fini radiciel}} \inserted{fini} sur un \unsure{sous-préschéma|sous-schéma} \unsure{fermé} \unsure{(ou ouvert)|(ouvert)} de $S$ ; \unsure{ce|le} morphisme est \deleted{un isomorphisme} \inserted{non ramifié} en \unsure{tt pt géom.}
% slanted
\margin{\uline{\unsure{ses pts|ex :} correspondant à une courbe non hyperelliptique}, \deleted{\illegible{3}} \uline{sa fibre} \deleted{\illegible{1}} au dessus d'un pt hyperelliptique \unsure{a un seul point} \deleted{\illegible{2}} \uline{géom.}, \unsure{isom. à} $\unsure{\mathrm{Spec}\,K[t]/t^2}$.}

\uline{Démonstration} \deleted{\unsure{1)} Commençons par représenter $F'(T)$, en commençant par remplacer le schéma $M$ sur $S$ Soit $T$ \unsure{faisons} le changement de base $T$} Considérons le foncteur $F'(T)$ analogue à $F(T)$, où \unsure{l'on} \deleted{nouveau foncteur représenté} met sur $C$ une structure supplémentaire : une section de $J^1_{C/S}$ sur $T$. \deleted{\unsure{D'après le numéro précédent|D'après les remarques générales} Il y a donc une \unsure{projection|morphisme} naturelle $F'(\deleted{\unsure{T'|T}}T)\to F(T)$, le produit fibré $\unsure{F'(T)\times_F|F(T)\times_F}$} $F'\times_F F'\mathrel{\unsure{\rightrightarrows|=}}F'$ étant \unsure{un} d'ailleurs, \unsure{si|d$^\circ$} $F'$ est représentable \inserted{par $M'$}, identifiable à \deleted{\unsure{un schéma}} $A\times M'$, \unsure{opérant sur|opérant en} $M'$ par des opérations de groupe.
% slanted
\margin{\deleted{\uline{\unsure{développer}} \uline{\unsure{en détail}}}}

% archive sheet number: 5

\scanpage{7}

\unsure{isons|isom|isom.}, on sait déjà qu'on a \deleted{\unsure{une}} un ouvert \uline{quasi-projectif} sur $S$ du schéma de Hilbert [car le \deleted{\unsure{pro}} degré de $C$ dans $A_T$ est connu \unsure{--|- -}]. \deleted{Il faut} La deuxième condition elle-même s'exprime par un ouvert plus petit.

On a ainsi prouvé que $F'$, \deleted{donc $F$}, est représentable quasi-projectivement, donc $F$ aussi. Nous voulons prouver que l'on a un \deleted{\uline{immersion}} \inserted{morphisme radiciel etc}. Par le critère \unsure{valuatif|habituel}, cela nous ramène aux \deleted{\unsure{de|deux cas}} particuliers suivants :

a) $S=\mathrm{Spec}(k)$ spectre d'un corps, alg. clos si on veut. a) \unsure{$\exists$}$M$ ne peut avoir que \deleted{\unsure{un}} \inserted{\unsure{au plus} deux} pt \unsure{ent $k$|sur $k$}, \inserted{conjugués par le groupe $\mathbb Z/2\subset\mathrm{Aut}_{\unsure{\mathrm{gr}|S}}(A)$} : c'est le \uline{théorème de Torelli} sous sa forme classique

b) \inserted{En \unsure{un|tout|fait} pt $x\in M$ autre que} $M\longrightarrow\mathrm{Spec}(k)$ est \uline{non ramifié}, i.e. \deleted{tt} \inserted{deux} pts de $M$ à valeurs ds $\mathrm{Spec}(k[\epsilon]/\epsilon^2)$ \deleted{\unsure{et ont pour}} \inserted{localisés en le même point de $M$} sont identiques. \deleted{\unsure{P|Pour}} Ce dernier point,

% archive sheet number: 6

\scanpage{8}

\deleted{En vertu du th. 1 \deleted{\illegible{1}}, il opère sans pt fixes, donc le quotient}

\uline{Si} $F'$ est représentable par \inserted{un schéma \unsure{qu.|quasi} projectif sur $S$} $M'$, alors $A$ opère sur $M'$, en vertu du th. 1 il opère sans pts fixes. \margin{à dégager} D'après un critère général \inserted{$A$ étant plat sur $S$}, il s'ensuit que $F$ est représentable \unsure{ssi|car} $M'/A$ existe ``raisonnablement'', et alors $M=M'/A$ représente $F$. Or on sait passer au quotient puisque $A$ est \uline{propre} et plat sur $S$,

Or $F'(T)$ s'explicite comme sous-ensemble de $\underline{\mathrm{Hilb}}_{A/S}(T)$, puisque les $C$ qu'on veut classifier \unsure{équivalent|équivalentes} aux $C$ \uline{plongés dans} $\deleted{\unsure{C_T}}A_T$, tels que le morphisme naturel $J^0_{C/S}\longrightarrow A_T$ déduit du \unsure{plongemt|plongement} soit un \uline{isom. de schémas Abéliens polarisés}, \deleted{\unsure{i.e. le degré est où la polarisation}} \deleted{Or, \unsure{le|la} exigeons que} Ceci implique d'ailleurs que, dans le morphisme $C\longrightarrow A_T$, \deleted{le sous-schéma fermé $C$ de $A_T$ soit} la polarisation induite sur $C$ soit de degré $g$, donc donné. D'après la théorie du schéma de Hilbert, si on fait abstraction du fait que $J^0_{C/S}\to A_T$ doive être un

% archive sheet number: 7

\scanpage{9}

signifie ceci : \deleted{Corollaire} \uline{Toute déformation \unsure{infinitésime|infinitésimale} du premier ordre de $C$ dans $A$} [si $C$ est plongé \unsure{jacobiennement|provisoirement|positionnée} ds $A$] \uline{est induit par une translation infinitésimale de $A$} (qui est bien déterminée d'ailleurs, puisque $C$ n'a pas d'automorphismes \inserted{infinitésimaux}).

Or \deleted{comme} les déformations infinitésimales en question sont définies par un élément de $H^0(C,N)$, où $N$ est le faisceau normal de $C$ dans $A$. Considérons la suite exacte
\diagram{\begin{tikzcd}[ampersand replacement=\&,column sep=small]
0\arrow[r,dash]\&\underline T\arrow[r]\&\underline V\arrow[r,dash]\&\underline N\arrow[r]\&0
\end{tikzcd}}
où $V$ est induit par le faisceau tangent à $A$ donc constant de rang $g$, $T$ est le faisceau tangent à $C$ donc d$^\circ$ inversible de degré $2(1-g)<0$, \deleted{et} \unsure{où|on} a donc $H^0(C,\underline T)=0$, $H^0(\unsure{C|O},\underline V)\simeq k^g$ $\underset{\unsure{\sim}}{=}$ \deleted{$\unsure{H^0(C}$} $\mathrm{Im}\,H^0(A,\unsure{\mathscr T|\mathscr J|\check{\Omega}}_{A/k})$, et la question est donc \unsure{exactement|seulement} si $H^0(C,\underline V)\mathbin{\unsure{-|\to}}H^0(\underline C,\underline N)$ est \uline{surjectif} i.e. $H^1(C,\underline T)\to H^1(C,\underline V)\xleftarrow{\sim}\deleted{\unsure{k}}H^1(A,\mathscr T_{A/k})$ injectif, i.e.

% archive sheet number: 8

\scanpage{10}

\deleted{\uline{Corollaire}} \uline{Toute déformation infinitésimale \inserted{du premier ordre} de la structure de $C$ qui induit une déformation infinitésimale triviale de la structure $J^0_{C/S}$, est triviale.}

\unsure{Transformant|Transposant|Traduisant} par dualité, cela équivaut à la question si
% slanted
\margin{dégager de façon \unsure{générale} \unsure{ttes} les compatibilités}
\[
H^0(C,\check V\otimes\underline\Omega)=H^0(C,\underline\Omega)\otimes H^0(C,\underline\Omega)\longrightarrow H^0(C,\underline\Omega^{\otimes 2})
\]
(où $\Omega=\Omega^1_{C/k}$) est \uline{surjectif}. Or ce dernier \uline{fait est vrai} \uline{ssi} $C$ \uline{est non hyperelliptique} ; dans le cas \deleted{\unsure{général}} hyperelliptique, on sait que le conoyau de l'hom précédent est de dim 1, donc le conoyau de $H^0(V)\mathbin{\unsure{-|\to}}H^0(N)$ est de dim 1, donc (s'agissant \unsure{du} point $x$ de $M$ \deleted{\unsure{issu}} sur $s\in S$) la fibre de $M$ en $s$ est $\mathrm{Spec}\,k[t]/t^2$.

(ii) On veut montrer que $M\mathbin{\unsure{-|\to}}S$ se factorise en un produit d'un morphisme \uline{propre} par une immersion. Par le critère valuatif, on est ramené à prouver ceci :

\uline{Corollaire} Supposons $S=\mathrm{Spec}(V)$, $V$ anneau de valuation discrète, \deleted{Soit} $\eta$ le pt générique de $S$, \deleted{\unsure{$C_\eta$ une}} On suppose $M\mathbin{\unsure{-|\to}}V$ surjectif i.e. que les fibres de $A$ sur $V$ sont des jacobiennes. \deleted{Alors} Soit $C_\eta$ une courbe sur

% archive sheet number: 9

\scanpage{11}

$\eta$ munie d'un isom \inserted{(compatible aux \unsure{pol$^s$|pol.})} de $J^0_{C_\eta/k(\eta)}\xrightarrow{\sim}A_\eta$. Alors il existe une courbe $C$ sur $S$, \deleted{tel que \unsure{sa réduction soit} $\unsure{C}$} dont la fibre générique soit $C_\eta$, et un isom de $J^0_{C/S}$ sur $A$.

En d'autres termes, la courbe $C_\eta$ ``se réduit bien'' [il s'ensuivra que $J^0_{C/S}$ est \unsure{néces$^{\rm mt}$|nécessairement} isomorphe à $A$ de façon unique de façon à induire l'isom donné sur la fibre générique].

\deleted{En fait, on \unsure{construit|construira} \unsure{directement|effectivement} $C$ dans $A$, car on peut supp. en tous les cas (pour ce qu'on veut en faire) que $V$ est \uline{complet} à corps résiduel alg. clos, donc $J^1_{C/S}$ a une section sur $S$}

\deleted{Soit} En fait, pour ce qu'on veut faire, on peut \deleted{supposer (quitte à dominer $V$ par un autre anneau de valuation) que} remplacer $M$ par $M'$, et par \unsure{suite} supposer déjà que $C_\eta$ correspond à un pt de $M'$ dans $k(\eta)$, i.e. que $\unsure{J^1_{C(\eta)/k(\eta)}|J_{C(\eta)/k(\eta)}|J_{C_\eta/k(\eta)}}$ est muni d'un pt rationnel sur $k(\eta)$, ou encore $C(\eta)$ \uline{plongé dans $A$}. \deleted{\unsure{par}} Soit $C$ la courbe sur $S$ adhérence de $C$. [\uline{Un résultat de Hoyt}

% archive sheet number: 10

\scanpage{12}

(en n'utilisant \unsure{pas|que} le fait que $\unsure{A_a|Aa}$ est une jacobienne) nous dit que \deleted{cette courbe} \inserted{$C_a$} a \unsure{ttes|bien} ses composantes de multiplicité $1$, et que $J^0_{C_a/k(a)}$ est isomorphe au produit des jacobiennes des composantes de $C_a$.]
% slanted
\margin{On n'a pas besoin de Hoyt, mais seulement de Matsusaka}

Mais un résultat de Matsusaka \inserted{(``géométrique'')} nous apprend que si on peut trouver sur $A$ un diviseur \inserted{\uline{positif}} \uline{irréductible} $\Theta_a$ tel que \deleted{$\deg(\Theta^n/n!)=1$, et que} $\frac{1}{n!}\deg\Theta_{a}^n=1$, \deleted{$\Theta^{n-1}/(n-1)!\ \sim$} et que $\Theta_a^{n-1}\equiv(n-1)!\,Z$ (équivalence numérique) où $Z$ est un cycle positif, alors $Z$ est \deleted{\unsure{ir}} une courbe non singulière irréductible, [et de plus -- mais nous n'en avons pas besoin -- que $A$ est isomorphe à sa jacobienne \deleted{\unsure{]}} par l'immersion $C\mathbin{\unsure{-|\to}}A$, enfin $\Theta$ définit la polarisation can. de la jacobienne].
% slanted
\margin{la première condition signifie que $\Theta$ définit une polarisation \uline{\unsure{unité}}}
\deleted{Or} \deleted{\unsure{Or} l'hypothèse \unsure{sur}} Considérons le diviseur $\Theta_\eta$ correspondant à $\unsure{J_\eta|\mathscr J_\eta}$, et son adhérence $\Theta$, d'où un $\Theta_a$. On \uline{sait} que $\deg(\Theta_\eta)^n/n!=1$ et $\Theta_\eta^{n-1}\equiv\deleted{\unsure{C_\eta}}(n-1)!\,C_\eta$, d'où $\deg(\Theta_a^n)/n!\unsure{\}|)}=1$ et $\Theta_a^{n-1}\equiv(n-1\unsure{\}|)}!\,C_a$. \deleted{\unsure{Comme}} \deleted{Si $C_a$ n'était irréductible, $\Theta_a$ ne serait pas}

% archive sheet number: 11

\scanpage{13}

\deleted{irréductible en vertu du th. de Matsusaka.} \deleted{\unsure{Alors}} \uline{On suppose déjà} que $A_{\unsure{a}}$ est une jacobienne \inserted{polarisée} ; \deleted{\unsure{alors qu'il existe|donc qu'il est} Or un résultat d'\uline{unicité} de Matsusaka affirme qu'il \unsure{n'}existe alors \inserted{à translation près} qu'un seul div. $D$ positif sur $A_a$ \deleted{\unsure{ayant les}}, \uline{irréductible ou non}, qui \deleted{ait les} satisfasse à deux équivalences numériques comme ci-dessus, et qu'alors} ce diviseur $\Theta_a$ est dans la classe de polarisation donc est \uline{irréductible} ; \deleted{donc} $C_a$ est irréductible et non \unsure{singulière|singulier} sur $k$. Donc $C$ se \deleted{réduit bien, et il} est simple sur $S$, et on achève maintenant \unsure{trivialement|facilement}.

La démonstration donne d'ailleurs ceci :

\deleted{Soit} \uline{Corollaire} 1 Soit $T$ \deleted{l'image} $=\overline{f(M)}$, où $f$ est le morphisme structural $M\to S$. Pour que $t\in T$ soit ds $f(M)$, il f. et s. que le diviseur de polarisation de $A_t$ [défini à translation près, en vertu de $\deg\Theta^g/g!=1$] soit \inserted{\unsure{ablt.|relt.|vraiment}} irréductible. \deleted{\unsure{§}}

% archive sheet number: 12

\scanpage{14}

\deleted{\uline{Corollaire 2} Soient $C,C'$ deux courbes sur $S$, munies d'isomorphismes de \deleted{schémas} schémas abéliens polarisés $\varphi:\deleted{\unsure{\varphi}}J^0_{C/S}\to A$ et $\varphi':J^0_{C'/S}\to A^{\unsure{*|\prime}}$. \deleted{Alors} On suppose que $S$ est réduit}

\uline{Corollaire 2} (i) Supposons que $S$ soit \inserted{\uline{réduit}, et} \uline{\unsure{connexe|normal}} \unsure{aux|avec} \deleted{$\unsure{s\in S}$} pts $s\in S$ tels que $A_{\deleted{\unsure{s}}s}$ soit une jacobienne \deleted{\unsure{non}} \uline{hyperelliptique}. Pour que $A$ soit un schéma jacobien \inserted{\uline{polarisé}}, il f. et s. que ses fibres le soient ; ou même que les fibres gén. (correspondant aux pts gén. des composantes de $S$) le soient, et que les diviseurs de polarisation des \inserted{\uline{\unsure{ttes|de}} les} fibres puissent être pris irréductibles
% slanted
\margin{il suffit que \deleted{\unsure{l'on}} $S$ soit réduit aux dits pts et que le morphisme $\unsure{\widetilde S|S'|s}\to S$ (\unsure{cf le n° suivant}) soit \uline{non ramifié} en $s$.}
% slanted
\margin{faux sous cette forme}

(ii) Supposons que $S$ soit \uline{réduit} aux \deleted{\unsure{pts}} $s\in S$ tels que $A_s$ soit hyperelliptique. Soient $C,C'$ deux courbes sur $S$, munies d'isom. de schémas abéliens polarisés $\varphi:J^0_{C/S}\mathbin{\unsure{-|\to}}A$, $\varphi':J^0_{C'/S}-A$. Alors il existe un $S$-isomorphisme $u:C-C'$ et un seul \deleted{tel que \unsure{qu}} compatible avec $\varphi$ et $\varphi'$.

% archive sheet number: 13
